\chapter{Statistiques}\label{ch:iv1:statistics}

\begin{ivtrains}
Évaluer ce qu'un échantillon nous apprend~: tests, intervalles et puissance, estimateurs et leurs erreurs, tests multiples, le bootstrap et les cas où il échoue, la normalité, et ce que vaut un ratio de Sharpe mesuré. La plupart des réponses se ramènent à une seule erreur type, calculée correctement.
\end{ivtrains}

\begin{question}[$\star$ \iqroles{QR, QT, DS, Analyste, Trader} \iqfirm{toutes} \iqlevel{junior}]\label{iq:iv1:statistics:1}
Une pièce donne 60 faces en 100 lancers. Testez si elle est équilibrée, au niveau de 5~\%, à l'aide de l'approximation normale. Le résultat est-il le même avec un calcul binomial exact~?
\end{question}

\begin{question}[$\star$ \iqroles{QR, QT, DS, Analyste} \iqfirm{toutes} \iqlevel{junior}]\label{iq:iv1:statistics:2}
Sur 400 jours, une stratégie affiche un rendement quotidien moyen de $0.1~\%$ et un écart-type quotidien de $1~\%$. Donnez un intervalle de confiance à 95~\% pour la vraie moyenne. Quelle taille d'échantillon rendrait l'intervalle deux fois moins large~?
\end{question}

\begin{question}[$\star$ \iqroles{QR, QT, DS, Analyste} \iqfirm{toutes} \iqlevel{junior}]\label{iq:iv1:statistics:3}
Un chercheur teste des idées dont 10~\% sont réellement bonnes. Le test a un taux de faux positifs de 5~\% et une puissance de 80~\%. Une idée vient de passer le test. Quelle est la probabilité que ce soit un fiasco~? Une p-valeur de $0.05$ signifie-t-elle qu'il y a 5~\% de chances que l'hypothèse nulle soit vraie~?
\end{question}

\begin{question}[$\star$ \iqroles{QR, QT, DS, Analyste} \iqfirm{toutes} \iqlevel{junior}]\label{iq:iv1:statistics:4}
Pourquoi la variance empirique divise-t-elle par $n - 1$ et non par $n$~? Montrez qu'elle rend l'estimateur sans biais. Que donne en moyenne l'estimateur naïf en $n$ pour des échantillons de taille 5~?
\end{question}

\begin{question}[$\star$ \iqroles{QR, QT, DS, Analyste} \iqfirm{toutes} \iqlevel{junior}]\label{iq:iv1:statistics:5}
Vous testez 20 signaux de trading indépendants, aucun ne fonctionnant, chacun au niveau de 5~\%. Quelle est la probabilité qu'un ou plusieurs passent le test~? Quel niveau par test maintiendrait cette probabilité à 5~\%~?
\end{question}

\begin{question}[$\star$ \iqroles{QR, QT, Analyste, Trader} \iqfirm{hedge fund} \iqlevel{junior}]\label{iq:iv1:statistics:6}
Mille gérants n'ont aucune compétence, et chacun bat le marché une année donnée avec une probabilité de 50~\%, indépendamment des autres. Après cinq ans, combien en moyenne ont battu le marché chaque année~? Que découvre un investisseur qui sélectionne un historique de cinq années sans accroc~?
\end{question}

\begin{question}[$\star$ \iqroles{QR, QT, DS, Analyste} \iqfirm{toutes} \iqlevel{junior}]\label{iq:iv1:statistics:7}
Une stratégie est positive 15 jours sur 20. Sans hypothèse sur la taille ni sur la forme des rendements, testez si son rendement quotidien médian est nul. Donnez la p-valeur exacte bilatérale.
\end{question}

\begin{question}[$\star$ \iqroles{QR, QT, DS} \iqfirm{toutes} \iqlevel{junior}]\label{iq:iv1:statistics:8}
L'algorithme A a produit les valeurs 5, 7 et 9, et l'algorithme B les valeurs 1, 2 et 3. Sans hypothèse sur les lois, testez si A donne des valeurs plus élevées, à l'aide d'un test de permutation. Quelle est la p-valeur unilatérale exacte~? Pourquoi aucun échantillon de cette taille ne peut-il descendre en dessous~?
\end{question}

\begin{question}[$\star\star$ \iqroles{QR, QT, DS, Analyste} \iqfirm{toutes} \iqlevel{junior}]\label{iq:iv1:statistics:9}
Un échantillon de 5 a pour moyenne 1 et pour écart-type 1. Donnez des intervalles à 95~\% pour la vraie moyenne, l'un avec la valeur normale $1.96$ et l'autre avec la valeur $t$ correcte, sachant que le point à 97.5~\% d'une loi $t$ à 4 degrés de liberté vaut $2.776$. Pourquoi diffèrent-ils, et quand la différence cesse-t-elle d'avoir de l'importance~?
\end{question}

\begin{question}[$\star\star$ \iqroles{QR, QT, DS, Analyste} \iqfirm{toutes} \iqlevel{junior}]\label{iq:iv1:statistics:10}
Deux séries ont une corrélation empirique de $0.5$ sur 103 points. Donnez la règle empirique pour l'erreur type d'une corrélation proche de zéro, puis un intervalle à 95~\% correct pour celle-ci, à l'aide de la transformation de Fisher. Pourquoi l'intervalle n'est-il pas symétrique autour de $0.5$~?
\end{question}

\begin{question}[$\star\star$ \iqroles{QR, QT, DS, Risque} \iqfirm{toutes} \iqlevel{junior}]\label{iq:iv1:statistics:11}
Sur 100 jours de trading, un système n'a connu aucune défaillance. Quelle borne supérieure pouvez-vous donner, à 95~\% de confiance, à sa probabilité de défaillance par jour~? Donnez le nombre exact et la règle empirique. Quelle est l'estimation naïve, et pourquoi n'est-elle pas utile~?
\end{question}

\begin{question}[$\star\star$ \iqroles{QR, QT, Analyste, Trader} \iqfirm{hedge fund} \iqlevel{junior}]\label{iq:iv1:statistics:12}
Une stratégie a un ratio de Sharpe annuel réel de $0.5$. Combien d'années de rendements quotidiens faut-il pour que son ratio de Sharpe mesuré soit, en moyenne, supérieur de deux erreurs types à zéro~? Qu'en est-il pour des ratios de Sharpe réels de 1 et de 2~? Qu'est-ce que cela dit de la longueur des historiques de performance~?
\end{question}

\begin{question}[$\star\star$ \iqroles{QR, QT, DS, Analyste} \iqfirm{hedge fund, teneur de marché} \iqlevel{junior}]\label{iq:iv1:statistics:13}
Un signal ajoute $0.05~\%$ par jour à un rendement dont l'écart-type quotidien est de $1~\%$. Pour un test bilatéral au niveau de 5~\% avec une puissance de 80~\%, combien de jours de données faut-il~? Donnez la formule et le résultat en années. Comment la réponse change-t-elle si l'effet est deux fois plus grand~?
\end{question}

\begin{question}[$\star\star$ \iqroles{QR, QT, DS} \iqfirm{toutes} \iqlevel{junior}]\label{iq:iv1:statistics:14}
Dix tests produisent les p-valeurs $0.001$, $0.008$, $0.012$, $0.041$, $0.042$, $0.06$, $0.074$, $0.205$, $0.212$ et $0.216$. Combien sont rejetées selon Bonferroni, avec un niveau global de $0.05$~? Combien selon Benjamini--Hochberg, avec un taux de fausses découvertes de $0.05$~? Que contrôle chacune des deux procédures~?
\end{question}

\begin{question}[$\star\star$ \iqroles{QR, QT, DS} \iqfirm{toutes} \iqlevel{junior}]\label{iq:iv1:statistics:15}
Les temps d'attente suivent une loi Gamma de forme $k$ et d'échelle $\theta$, dont la moyenne vaut $k\theta$ et la variance $k\theta^2$. Un échantillon a pour moyenne 4 et pour variance 8. Donnez les estimations par la méthode des moments. Pourquoi préféreriez-vous malgré tout les estimations du maximum de vraisemblance, si vous savez les calculer~?
\end{question}

\begin{question}[$\star\star$ \iqroles{QR, QT, DS} \iqfirm{toutes} \iqlevel{junior}]\label{iq:iv1:statistics:16}
Vous disposez de $n$ temps d'attente exponentiels indépendants de taux $\lambda$. Dérivez l'estimateur du maximum de vraisemblance du taux. Est-il sans biais~? Donnez la version sans biais, et sa valeur pour $n = 5$ lorsque la moyenne empirique vaut 2.
\end{question}

\begin{question}[$\star\star$ \iqroles{QR, QT, DS} \iqfirm{toutes} \iqlevel{junior}]\label{iq:iv1:statistics:17}
Les comptages quotidiens suivent une loi de Poisson de moyenne $\lambda$, et vous disposez de $n$ jours. Montrez que la moyenne empirique atteint la borne inférieure sur la variance d'un estimateur sans biais. Pour estimer la probabilité d'un jour sans aucun événement, comparez la fraction de jours nuls avec $e^{-\bar{X}}$, pour $\lambda = 2$, en variance asymptotique. Lequel est le meilleur, et que suppose le meilleur des deux~?
\end{question}

\begin{question}[$\star\star$ \iqroles{QR, QT, DS, QD} \iqfirm{toutes} \iqlevel{junior}]\label{iq:iv1:statistics:18}
Vous tirez un échantillon bootstrap de taille $n$ avec remise à partir de $n$ points de données. Quelle fraction des points d'origine contient-il en moyenne, pour $n$ grand~? Que estime le bootstrap, et pourquoi cela fonctionne-t-il, en termes simples~?
\end{question}

\begin{question}[$\star\star$ \iqroles{QR, QT, DS, Risque} \iqfirm{toutes} \iqlevel{junior}]\label{iq:iv1:statistics:19}
Les rendements quotidiens suivent un AR(1) de coefficient $\rho = 0.5$. De quel facteur la variance de la moyenne empirique dépasse-t-elle la valeur calculée en supposant l'indépendance~? Que fournit un bootstrap simple, et que faut-il faire à la place~?
\end{question}

\begin{question}[$\star\star$ \iqroles{QR, QT, DS, Risque} \iqfirm{toutes} \iqlevel{junior}]\label{iq:iv1:statistics:20}
Mille rendements quotidiens ont une asymétrie de $-0.5$ et un kurtosis en excès de 3. Donnez les erreurs types de ces deux nombres sous l'hypothèse de normalité. Calculez la statistique de Jarque--Bera, et dites à quoi on la compare. Quelle est la faiblesse de ce test pour les données financières~?
\end{question}

\begin{question}[$\star\star$ \iqroles{QR, QT, Risque, Trader, Analyste} \iqfirm{toutes} \iqlevel{junior}]\label{iq:iv1:statistics:21}
Un rendement quotidien a un écart-type de $1~\%$, et aujourd'hui il vaut $-5~\%$. Si les rendements sont normaux, à quelle fréquence devrait-on observer un tel jour, en années~? À quelle fréquence si les rendements suivent une loi $t$ de Student à 3 degrés de liberté, remise à l'échelle pour le même écart-type~? Que concluez-vous~?
\end{question}

\begin{question}[$\star\star$ \iqroles{QR, QT, DS} \iqfirm{toutes} \iqlevel{junior}]\label{iq:iv1:statistics:22}
Un dé est lancé 120 fois, et les faces 1 à 6 sortent respectivement 30, 18, 20, 22, 16 et 14 fois. Effectuez un test du khi-deux d'équilibre, sachant que le point à 95~\% à 5 degrés de liberté vaut $11.07$. Un rapide coup d'œil suggère que la face 1 est favorisée. Testeriez-vous cette seule face, et comment~?
\end{question}

\begin{question}[$\star\star$ \iqroles{QR, QT, DS, Trader} \iqfirm{toutes} \iqlevel{junior}]\label{iq:iv1:statistics:23}
Une suite de 100 jours à la hausse et à la baisse contient 30 séries, une série étant un bloc maximal de même signe. Si les jours sont indépendants et équilibrés, quelles sont la moyenne et l'écart-type du nombre de séries~? Que suggère la valeur 30~? Que suggérerait la valeur 65~?
\end{question}

\begin{question}[$\star\star$ \iqroles{QR, QT, DS, LLE} \iqfirm{teneur de marché, hedge fund} \iqlevel{junior}]\label{iq:iv1:statistics:24}
Deux algorithmes d'exécution traitent chacun les mêmes 100 ordres. Leurs glissements ont chacun un écart-type de 10 points de base, une corrélation de $0.9$ entre les ordres, et une différence moyenne de 1 point de base. Quel est le score $z$ si on les compare comme deux échantillons indépendants~? Et si on les compare par paires~? Lequel est correct, et que cela suggère-t-il pour la conception d'un test~?
\end{question}

\begin{question}[$\star\star$ \iqroles{QR, QT, DS, Analyste} \iqfirm{toutes} \iqlevel{junior}]\label{iq:iv1:statistics:25}
Deux estimations indépendantes valent 0 et 3, chacune avec une erreur type de 1. Leurs intervalles à 95~\% se chevauchent. La différence est-elle significative au niveau de 5~\%~? Trouvez la règle générale qui relie le chevauchement à la significativité.
\end{question}

\begin{question}[$\star\star$ \iqroles{QR, QT, DS, Risque} \iqfirm{toutes} \iqlevel{junior}]\label{iq:iv1:statistics:26}
Pour dix points, $x$ va de 1 à 10, et $y = x$, sauf que le dernier $y$ vaut 100. Calculez les corrélations de Pearson et de Spearman. Laquelle donne le résumé le plus juste, et dans quel cas vous méfieriez-vous de chacune d'elles~?
\end{question}

\begin{question}[$\star\star\star$ \iqroles{QR, QT, DS} \iqfirm{hedge fund} \iqlevel{junior, confirmé}]\label{iq:iv1:statistics:27}
Vous estimez la variance d'une loi normale à partir de $n$ tirages indépendants, en utilisant $\frac{1}{c}\sum (X_i - \bar{X})^2$. Trouvez le diviseur $c$ qui donne l'erreur quadratique moyenne la plus petite. Pour $n = 10$, donnez l'erreur quadratique moyenne, en unités de $\sigma^4$, pour $c = n - 1$, pour $n$ et pour le meilleur $c$. Comment cela se concilie-t-il avec le fait que le caractère sans biais est une qualité~?
\end{question}

\begin{question}[$\star\star\star$ \iqroles{QR, QT, DS} \iqfirm{hedge fund} \iqlevel{junior, confirmé}]\label{iq:iv1:statistics:28}
Vous devez choisir entre deux hypothèses exactes à partir de $n$ tirages indépendants~: la loi est $N(0, 1)$, ou bien elle est $N(1, 1)$. Quel est le meilleur test, et pourquoi~? Quelle doit être la taille $n$ pour que les deux taux d'erreur soient au plus égaux à 5~\%, s'ils sont fixés égaux~?
\end{question}

\begin{question}[$\star\star\star$ \iqroles{QR, QT, DS, Analyste} \iqfirm{hedge fund, teneur de marché} \iqlevel{junior, confirmé}]\label{iq:iv1:statistics:29}
Vous menez une expérience sur un effet nul. Après chacun des 10 lots équidistants, vous calculez un score $z$ sur l'ensemble des données disponibles, et vous arrêtez dès que $|z| > 1.96$. Quel est à peu près votre taux réel de faux positifs~? Que se passe-t-il sans limite sur le nombre de regards~? Quel seuil par regard rétablirait 5~\%, et comment appelle-t-on la méthode~?
\end{question}

\begin{question}[$\star\star\star$ \iqroles{QR, QT, DS} \iqfirm{toutes} \iqlevel{junior, confirmé}]\label{iq:iv1:statistics:30}
Vous tirez $n$ valeurs indépendantes d'une loi uniforme sur $[0, \theta]$, et vous voulez un intervalle de confiance pour $\theta$ à partir du maximum de l'échantillon. Montrez qu'un bootstrap naïf du maximum échoue, en calculant la probabilité que le maximum bootstrap soit égal au maximum de l'échantillon, pour $n$ grand. À quoi ressemble la vraie loi d'échantillonnage, et quelle méthode utiliseriez-vous~?
\end{question}

\begin{question}[$\star\star\star$ \iqroles{QR, QT, Analyste, Risque} \iqfirm{hedge fund} \iqlevel{junior, confirmé}]\label{iq:iv1:statistics:31}
Une stratégie présente un ratio de Sharpe mensuel de $s = 0.3$ mesuré sur $n = 120$ mois. Déduisez l'erreur type de l'estimation sous l'hypothèse de normalité. De combien l'intervalle s'élargit-il si les rendements mensuels ont une asymétrie de $-1$ et un kurtosis de 6, sachant que la variance devient $\bigl(1 - \gamma_3 s + \frac{\gamma_4 - 1}{4}s^2\bigr)/n$~? Pourquoi s'élargit-il~?
\end{question}

\begin{question}[$\star\star\star$ \iqroles{QR, QT, Analyste, Risque} \iqfirm{hedge fund} \iqlevel{junior, confirmé}]\label{iq:iv1:statistics:32}
Les rendements mensuels ont un ratio de Sharpe de $0.3$ et une autocorrélation d'ordre 1 de $0.2$, suivant un AR(1). La règle habituelle multiplie par $\sqrt{12}$. Quel est le ratio de Sharpe annuel correct~? Que se passe-t-il avec une autocorrélation négative~? Quels types de stratégies affichent une autocorrélation positive dans leurs rendements publiés, et pourquoi~?
\end{question}

\begin{question}[$\star\star\star$ \iqroles{QR, QT, DS} \iqfirm{hedge fund} \iqlevel{junior, confirmé}]\label{iq:iv1:statistics:33}
Deux stratégies ont le même ratio de Sharpe annuel vrai et une corrélation de $0.8$, et vous disposez de 5 ans de rendements quotidiens pour chacune. Quelle est l'erreur type de la différence des deux ratios de Sharpe mesurés~? Comment se compare-t-elle au calcul qui les traite comme indépendantes~? Qu'est-ce que cela implique pour le choix entre des stratégies similaires~?
\end{question}

\begin{question}[$\star\star\star$ \iqroles{QR, QT, Risque, DS} \iqfirm{banque, hedge fund} \iqlevel{junior, confirmé}]\label{iq:iv1:statistics:34}
Vous estimez le quantile à 1~\% d'un rendement quotidien, la valeur en risque (VaR), à partir de 1000 jours de rendements normaux. Déduisez l'erreur type du quantile empirique, et comparez-la à l'erreur de $2.326\hat{\sigma}$, en unités de $\sigma$. Pourquoi la première est-elle plus grande~? Quand faudrait-il tout de même choisir la première~?
\end{question}

\begin{question}[$\star\star\star$ \iqroles{QR, QT, Risque, DS} \iqfirm{toutes} \iqlevel{junior, confirmé}]\label{iq:iv1:statistics:35}
Un échantillon de 101 rendements a un écart-type de $2~\%$. Quelle est l'erreur type de cette estimation si les rendements sont normaux~? S'ils ont un kurtosis de 6~? Déduisez la règle. Qu'est-ce que cela change à l'intervalle du chi-deux habituel~?
\end{question}

\begin{question}[$\star\star\star$ \iqroles{QR, QT, DS} \iqfirm{hedge fund} \iqlevel{junior, confirmé}]\label{iq:iv1:statistics:36}
Vous mesurez une corrélation de $0.02$ entre un signal et les rendements du lendemain. La moitié de la variance du signal est du bruit de mesure, et les trois quarts de la variance du rendement du jour sont un bruit sans rapport avec ce que vous pourriez prévoir. Quelle est la corrélation entre le vrai signal et la partie prévisible des rendements~? Établissez la formule. Qu'est-ce que cela dit des corrélations brutes en finance~?
\end{question}

\begin{question}[$\star\star\star$ \iqroles{QR, QT, DS, Risque} \iqfirm{hedge fund, teneur de marché} \iqlevel{junior, confirmé}]\label{iq:iv1:statistics:37}
Le rendement d'un actif a une volatilité annuelle de 20~\%, et vous disposez de 10 ans de données quotidiennes. Calculez les erreurs types de la moyenne annuelle et de la volatilité annuelle, en points de pourcentage. Comment chacune change-t-elle si vous échantillonnez au mois, ou à la minute~? Quelle conclusion pratique en découle, et qu'est-ce qui vous empêche de pousser sans limite le résultat sur la volatilité~?
\end{question}

\begin{question}[$\star\star\star\star$ \iqroles{QR, QT, DS} \iqfirm{hedge fund} \iqlevel{confirmé}]\label{iq:iv1:statistics:38}
Vous observez $X \sim N(\theta, I_d)$ en dimension $d$ et vous estimez la moyenne $\theta$ avec pour perte l'erreur quadratique totale. Montrez que l'estimateur de James--Stein $\bigl(1 - (d - 2)/|X|^2\bigr)X$ fait mieux que $X$ lui-même pour tout $\theta$ lorsque $d \ge 3$. Donnez son risque en $\theta = 0$ pour $d = 10$. Pourquoi est-ce étrange, et comment l'utiliseriez-vous sur un ensemble de rendements moyens d'actifs~?
\end{question}

\begin{question}[$\star\star\star\star$ \iqroles{QR, QT, DS} \iqfirm{hedge fund} \iqlevel{confirmé}]\label{iq:iv1:statistics:39}
Un échantillon de taille $n$ provient d'une loi de Cauchy de centre $\theta$, de densité $\frac{1}{\pi(1 + (x - \theta)^2)}$. Calculez l'information de Fisher par observation. Comparez les variances asymptotiques de l'estimateur du maximum de vraisemblance, de la médiane empirique et de la moyenne empirique. Comment calculeriez-vous le premier en pratique, sachant que sa vraisemblance peut présenter plusieurs maxima locaux~?
\end{question}

\begin{question}[$\star\star\star\star$ \iqroles{QR, QT, DS} \iqfirm{hedge fund} \iqlevel{confirmé}]\label{iq:iv1:statistics:40}
Prouvez que la procédure de Benjamini--Hochberg au niveau $\alpha$ contrôle le taux de fausses découvertes à $\alpha m_0/m$ lorsque les $m$ p-valeurs sont indépendantes, avec $m_0$ vraies hypothèses nulles dont les p-valeurs sont uniformes. Indication~: pour une seule vraie hypothèse nulle, demandez-vous quel serait le nombre de rejets si sa p-valeur était fixée à 0.
\end{question}
